% GENERATED FILE. Edit canonical lessons, then run npm run generate:books.
% canonical-source-hash: fnv1a64:4d69533b1fd8043c
% canonical-lessons: KA-C85-slowly, KA-C85-oncemore, KA-C85-loudly, KA-C85-meaning, KA-C85-writeit

\chapter{Slowly, Once More}
\label{ch:ka-slowly-once-more}
\hlchaptermodality{85}

\begin{chapteropening}
\textbf{By the end of this chapter:} \emph{I can ask someone to speak slowly, say it once more, speak louder, tell me the meaning, and write it down.}

Complete the last lesson of chapter 85: i can ask someone to speak slowly, say it once more, speak louder, tell me the meaning, and write it down.
\end{chapteropening}

\section[nidhānavāgi]{\kn{ನಿಧಾನವಾಗಿ} (nidhānavāgi) — slowly}

\label{lesson:KA-C85-slowly}



\begin{quote}
Before the new one: say the Kannada for hospital, then the Kannada for a tablet.
\end{quote}

\subsection*{You'll want to know: \kn{ನಿಧಾನವಾಗಿ}}
\textbf{\kn{ನಿಧಾನವಾಗಿ}} (\emph{nidhānavāgi}) — \textquotedblleft{}slowly\textquotedblright{}. The ending \textbf{‑\kn{ಆಗಿ}} (\emph{-āgi}) turns a word into a manner, \textquotedblleft{}in a slow way\textquotedblright{}, like \textbf{\kn{ಚೆನ್ನಾಗಿ}} (\emph{cennāgi}), \textquotedblleft{}well\textquotedblright{}.

\textbf{\kn{ನಿಧಾನವಾಗಿ} \kn{ಮಾತಾಡಿ}} (\emph{nidhānavāgi mātāḍi}) — \textquotedblleft{}please speak slowly\textquotedblright{}.

\begin{grammarlens}[title={What you've built}]
The most useful request a learner can make.
\end{grammarlens}

\subsection*{Your turn}
\begin{itemize}
\raggedright
  \item \emph{Say it:} \emph{nidhānavāgi}
  \item \emph{Say it:} \emph{nidhānavāgi}, once more
  \item \emph{Say it:} say \emph{nidhānavāgi mātāḍi}
  \item \emph{Recall:} say the Kannada for hospital, then the Kannada for a tablet, then say \emph{nidhānavāgi} again
\end{itemize}

\begin{tcolorbox}[breakable,colback=teal!4,colframe=teal!35!black,title={Before you move on}]
What does \kn{ನಿಧಾನವಾಗಿ} mean? (\textquotedblleft{}Slowly\textquotedblright{}.) Say it once more.
\end{tcolorbox}
\section[mattomme]{\kn{ಮತ್ತೊಮ್ಮೆ} (mattomme) — once more}

\label{lesson:KA-C85-oncemore}



\begin{quote}
Before the new one: say the Kannada for a tablet, then the Kannada for slowly.
\end{quote}

\subsection*{You'll want to know: \kn{ಮತ್ತೊಮ್ಮೆ}}
\textbf{\kn{ಮತ್ತೊಮ್ಮೆ}} (\emph{mattomme}) — \textquotedblleft{}once more, one more time\textquotedblright{}.

\textbf{\kn{ಮತ್ತೊಮ್ಮೆ} \kn{ಹೇಳಿ}} (\emph{mattomme hēḷi}) — \textquotedblleft{}please say it once more\textquotedblright{}, with \textbf{\kn{ಹೇಳಿ}}, \textquotedblleft{}please tell\textquotedblright{}, which you already know.

\begin{grammarlens}[title={What you've built}]
The request that comes after slowly.
\end{grammarlens}

\subsection*{Your turn}
\begin{itemize}
\raggedright
  \item \emph{Say it:} \emph{mattomme}
  \item \emph{Say it:} \emph{mattomme}, once more
  \item \emph{Say it:} say \emph{mattomme hēḷi}
  \item \emph{Recall:} say the Kannada for a tablet, then the Kannada for slowly, then say \emph{mattomme} again
\end{itemize}

\begin{tcolorbox}[breakable,colback=teal!4,colframe=teal!35!black,title={Before you move on}]
What does \kn{ಮತ್ತೊಮ್ಮೆ} mean? (\textquotedblleft{}Once more\textquotedblright{}.) Say it once more.
\end{tcolorbox}
\section[jōrāgi]{\kn{ಜೋರಾಗಿ} (jōrāgi) — loudly}

\label{lesson:KA-C85-loudly}



\begin{quote}
Before the new one: say the Kannada for slowly, then the Kannada for once more.
\end{quote}

\subsection*{You'll want to know: \kn{ಜೋರಾಗಿ}}
\textbf{\kn{ಜೋರಾಗಿ}} (\emph{jōrāgi}) — \textquotedblleft{}loudly\textquotedblright{}. It has the same \textbf{‑\kn{ಆಗಿ}} ending as \textbf{\kn{ನಿಧಾನವಾಗಿ}}.

\textbf{\kn{ಸ್ವಲ್ಪ} \kn{ಜೋರಾಗಿ}} (\emph{svalpa jōrāgi}) — \textquotedblleft{}a little louder\textquotedblright{}, with \textbf{\kn{ಸ್ವಲ್ಪ}}, \textquotedblleft{}a little\textquotedblright{}.

\begin{grammarlens}[title={What you've built}]
A second way of saying how.
\end{grammarlens}

\subsection*{Your turn}
\begin{itemize}
\raggedright
  \item \emph{Say it:} \emph{jōrāgi}
  \item \emph{Say it:} \emph{jōrāgi}, once more
  \item \emph{Say it:} say \emph{svalpa jōrāgi}
  \item \emph{Recall:} say the Kannada for slowly, then the Kannada for once more, then say \emph{jōrāgi} again
\end{itemize}

\begin{tcolorbox}[breakable,colback=teal!4,colframe=teal!35!black,title={Before you move on}]
What does \kn{ಜೋರಾಗಿ} mean? (\textquotedblleft{}Loudly\textquotedblright{}.) Say it once more.
\end{tcolorbox}
\section[artha]{\kn{ಅರ್ಥ} (artha) — meaning}

\label{lesson:KA-C85-meaning}



\begin{quote}
Before the new one: say the Kannada for once more, then the Kannada for loudly.
\end{quote}

\subsection*{You'll want to know: \kn{ಅರ್ಥ}}
\textbf{\kn{ಅರ್ಥ}} (\emph{artha}) — \textquotedblleft{}meaning\textquotedblright{}. You have seen it inside \textbf{\kn{ಅರ್ಥಮಾಡಿಕೊ}}, \textquotedblleft{}to understand\textquotedblright{}, which is \textquotedblleft{}to make the meaning one's own\textquotedblright{}.

\textbf{\kn{ಇದರ} \kn{ಅರ್ಥ} \kn{ಏನು}?} (\emph{idara artha ēnu?}) — \textquotedblleft{}what is the meaning of this?\textquotedblright{}

\begin{grammarlens}[title={What you've built}]
The piece hidden inside the verb to understand.
\end{grammarlens}

\subsection*{Your turn}
\begin{itemize}
\raggedright
  \item \emph{Say it:} \emph{artha}
  \item \emph{Say it:} \emph{artha}, once more
  \item \emph{Say it:} ask \emph{idara artha ēnu?}
  \item \emph{Recall:} say the Kannada for once more, then the Kannada for loudly, then say \emph{artha} again
\end{itemize}

\begin{tcolorbox}[breakable,colback=teal!4,colframe=teal!35!black,title={Before you move on}]
What does \kn{ಅರ್ಥ} mean? (\textquotedblleft{}Meaning\textquotedblright{}.) Say it once more.
\end{tcolorbox}
\section[baredu koḍi]{\kn{ಬರೆದು} \kn{ಕೊಡಿ} (baredu koḍi) — please write it down}

\label{lesson:KA-C85-writeit}



\begin{quote}
Before the new one: say the Kannada for loudly, then the Kannada for meaning.
\end{quote}

\subsection*{You'll want to know: \kn{ಬರೆದು} \kn{ಕೊಡಿ}}
\textbf{\kn{ಬರೆದು} \kn{ಕೊಡಿ}} (\emph{baredu koḍi}) — \textquotedblleft{}please write it down (for me)\textquotedblright{}, literally \textquotedblleft{}having written, give\textquotedblright{}. Both verbs are ones you know: \textbf{\kn{ಬರೆ}}, \textquotedblleft{}write\textquotedblright{}, and \textbf{\kn{ಕೊಡು}}, \textquotedblleft{}give\textquotedblright{}.

When a word will not stay in the ear, ask for it on paper.

\begin{grammarlens}[title={What you've built}]
Two known verbs make one request.
\end{grammarlens}

\subsection*{Your turn}
\begin{itemize}
\raggedright
  \item \emph{Say it:} \emph{baredu koḍi}
  \item \emph{Say it:} \emph{baredu koḍi}, once more
  \item \emph{Say it:} say \emph{baredu koḍi}
  \item \emph{Recall:} say the Kannada for loudly, then the Kannada for meaning, then say \emph{baredu koḍi} again
\end{itemize}

\begin{tcolorbox}[breakable,colback=teal!4,colframe=teal!35!black,title={Before you move on}]
What does \kn{ಬರೆದು} \kn{ಕೊಡಿ} mean? (\textquotedblleft{}Please write it down\textquotedblright{}.) Say it once more.
\end{tcolorbox}
